
This work relates to three lines of research: studies of benchmark generalization and distribution shift, analyses of benchmark concentration and selection bias, and methods for measuring ranking stability.

\subsubsection{Benchmark Generalization and Distribution Shift}

Recht et al. (2019) created ImageNet-V2 by replicating the original ImageNet data collection methodology. All tested models experienced 11-14\% accuracy drops on the new test set. They observed that ''accuracy gains on the original test set translate to roughly the same gain on the new test set,'' suggesting proportional degradation. The present work quantifies this observation via ranking correlation.

Taori et al. (2020) introduced effective robustness, showing a linear relationship between in-distribution accuracy and out-of-distribution accuracy across multiple distribution shifts. A linear relationship implies that rankings should be preserved, which is consistent with the findings reported here. However, Taori et al. did not explicitly compute ranking correlations.

Miller et al. (2021) studied distribution shift in question answering, finding that model rankings can shift substantially when evaluation data changes. Their finding of ranking instability in NLP contrasts with the stability found here for vision, suggesting domain-specific effects.

\subsubsection{Benchmark Concentration and Selection Bias}

Koch et al. (2021) documented that the top 10\% of NLP datasets account for the same usage as the remaining 90\% combined. This concentration creates conditions for benchmark-specific optimization.

Dehghani et al. (2021) demonstrated the benchmark lottery phenomenon: model rankings on SuperGLUE tasks depend heavily on which tasks are selected. By re-computing aggregate scores with different task combinations, they showed that apparent leaders may be artifacts of task selection. While their work focused on task selection within a benchmark, the present work focuses on generalization across benchmark variants.

Beyer et al. (2020) documented ceiling effects and label noise issues with ImageNet, questioning whether continued progress reflects genuine capability improvements. These concerns motivate investigation of whether ImageNet-based rankings transfer to cleaner evaluation settings.

\subsubsection{Ranking Stability Measurement}

Kendall-$\tau$ and Spearman-$\rho$ are standard metrics for comparing rankings across contexts. Prior work in information retrieval has used these metrics to assess ranking stability under query variations \citep{voorhees2000variations}. In machine learning evaluation, these metrics are less commonly applied. The present work fills this gap for benchmark comparison.

